
% =====================================================================
\section{The order parameter and the ground state}
\label{sec:order}
% =====================================================================

Section~\ref{sec:cjt} showed that the four vector pairs along the extra axes have the strongest
pairing. This section determines how they combine: the structure of the order parameter, the energy
as a function of it, the ground state, and the place of the hidden-axis singlet in it. It completes
objective O4.

\subsection{The gap matrix}

The condensate is described by the expectation values of the four vector pair operators along the
extra axes,
\begin{equation}
\Delta^a\propto\big\langle\psi^T\Cm\gamma^a\psi\big\rangle,\qquad a\in\{0,5,6,7\},
\label{eq:Deltaa}
\end{equation}
four complex numbers. On the eight internal states of the pair $(\vec p,-\vec p)$ they define the
\emph{gap matrix}
\begin{equation}
\MM(\vec p)=\sum_{a}\Delta^a\,\frac{\Phi_a(\vec p)}{\sqrt{N_a/8}},\qquad\Phi_a=U(\vec p)^T\Cm\gamma^aU(-\vec p),
\label{eq:gapmatrix}
\end{equation}
with the two-particle contents $\Phi_a$ of Eq.~(\ref{eq:content}) and their norms $N_a$. The energies
of the quasiparticles are $\sqrt{\xi^2+d_r^2}$, where $d_r^2$, $r=1,\dots,8$, are the eigenvalues of
$\MM\MM^\dagger$. A state is called \emph{unitary} if $\MM\MM^\dagger$ is a multiple of the unit
matrix, so that all eight internal states have the same gap \cite{Leggett1975}.

\subsection{The energy as a function of the order parameter}
\label{sec:GL}

Near the critical temperature the free energy can be expanded in the order parameter. At weak
coupling the expansion of Ginzburg and Landau, as derived from the microscopic theory by Gor'kov
\cite{Gorkov1959}, is
\begin{align}
\mathcal F={}&a(T)\,\Tr\big[\MM\MM^\dagger\big]+\frac b2\,\Tr\big[(\MM\MM^\dagger)^2\big]\nonumber\\
&+K_{\rm GL}\,\Tr\big[\nabla\MM\cdot\nabla\MM^\dagger\big]+\dots,\nonumber\\
a(T)={}&\frac{N(0)}2\ln\frac{T}{T_c},\qquad b=\frac{7\zeta(3)}{16\pi^2T_c^2}N(0),\nonumber\\
K_{\rm GL}={}&\frac{7\zeta(3)}{96\pi^2T_c^2}N(0)\,v_F^2,
\label{eq:GL}
\end{align}
with the Fermi velocity $v_F$ and the density of states $N(0)$ of Eq.~(\ref{eq:lambda_coupling}). For
a unitary state this is the textbook energy of four independent copies of a spin-$\tfrac12$
superconductor. Section~\ref{sec:checks} confirmed the three coefficients from the CJT potential.

The form of the quartic term is decided by one matrix identity. The program verifies it from the
matrices of both spacetimes, with a largest deviation of {{GEN:gl_dev_@SIG@}} for the present one:
\begin{align}
\frac18\Tr\big[\MM\MM^\dagger\big]&=\sum_a|\Delta^a|^2\equiv|\Delta|^2,\nonumber\\
\frac18\Tr\big[(\MM\MM^\dagger)^2\big]&=2\,|\Delta|^4-\big|\eta_{ab}\,\Delta^a\Delta^b\big|^2 ,
\label{eq:quartic}
\end{align}
where $\eta_{ab}$ is the metric of the spacetime restricted to the four extra axes,
\begin{equation}
\eta_{ab}=\SIG{\mathrm{diag}(+1,-1,-1,-1)}{\mathrm{diag}(-1,-1,-1,-1)}\qquad(a,b\in\{0,5,6,7\})
\label{eq:etaab}
\end{equation}
for the \SIG{$4+4$}{$7+1$} spacetime: \SIG{the hidden axis is space-like and the other three are
time-like. On the $7+1$ spacetime all four are space-like, and $\eta_{ab}$ is proportional to the
unit matrix.}{all four extra axes are space-like. On the $4+4$ spacetime the hidden axis is
space-like and the other three are time-like, and $\eta_{ab}=\mathrm{diag}(+1,-1,-1,-1)$.} \emph{The
signature of the spacetime thus enters the energy of the superfluid through the quadratic form
$\eta_{ab}\Delta^a\Delta^b$.} It is the invariant form (\ref{eq:invform}) of the internal symmetry.

Equation (\ref{eq:quartic}) is a statement about the sum of the fourth powers of the gaps. The
complete statement, which the program also verifies, is that the eight eigenvalues of
$\MM\MM^\dagger$ form two fourfold sets,
\begin{equation}
d_\pm^2=|\Delta|^2\pm\sqrt{|\Delta|^4-\big|\eta_{ab}\Delta^a\Delta^b\big|^2}.
\label{eq:twogaps}
\end{equation}
A state is unitary, $d_+=d_-$, exactly when $|\eta_{ab}\Delta^a\Delta^b|=|\Delta|^2$.

\subsection{The ground state, and the hidden-axis singlet}
\label{sec:ground}

At fixed $|\Delta|^2$ the quartic term (\ref{eq:quartic}) is smallest when
$|\eta_{ab}\Delta^a\Delta^b|$ is largest, that is, equal to $|\Delta|^2$. The ground states are
therefore the unitary states: all eight quasiparticle branches have the same gap, and the
condensation energy is the largest possible. The minimizing configurations are
\begin{equation}
\SIG{\Delta^0=|\Delta|\,e^{i\varphi}\,n^0,\qquad\Delta^t=\pm i\,|\Delta|\,e^{i\varphi}\,n^t\quad(t=5,6,7)}{\Delta^a=|\Delta|\,e^{i\varphi}\,n^a}
\label{eq:groundstates}
\end{equation}
with a phase $\varphi$ and a \emph{real} unit vector
\SIG{$(n^0,n^5,n^6,n^7)$: the components along the time-like axes carry a relative phase $\pm i$.}{$n^a$
in the space of the four extra axes.} The \emph{hidden-axis singlet} of hypothesis H1,
\begin{equation}
\Delta^a=\Delta\,\delta^{a0},\qquad\big\langle\psi_{\tg}^T\,\Cm_{\tg}\gamma_{\tg}^0\,\psi_{\tg}\big\rangle\neq0,
\label{eq:singlet}
\end{equation}
is a member of this family. Its properties were proved in Sec.~\ref{sec:hiddenpair}: the pair has
ordinary spin 0; it is invariant under the rotations of ordinary space and of the coordinates
$x5,x6,x7$; and it is a scalar under the Lorentz group of the seven coordinates other than $x0$.

\emph{How the hidden axis is singled out.} By Sec.~\ref{sec:internal} the order parameter is a
four-vector of the internal $SO(4)$ of the zero-mode sector, and all directions of this four-vector
have exactly the same energy.
%%<1671>
On the $7+1$ spacetime the direction $n^a$ is chosen spontaneously, and calling the chosen one ``the
hidden axis'' is a choice of coordinates. The symmetry breaking is $SO(4)\to SO(3)$ together with
$U(1)_F\to Z_2$. Besides the sound mode of the superfluid, the Goldstone mode of the fermion number,
there are three Goldstone modes, the rotations of $n^a$. Because the $SO(4)$ is a symmetry of the full
theory, nothing within the flat $7+1$ model prefers one direction. In the second stage of the program
the hidden coordinate is curled up and the three others are treated differently; this would break the
$SO(4)$ explicitly and decide the direction. It is not implemented here.
%%</1671>
%%<1644>
On the $4+4$ spacetime the components that form the real four-vector are
$(\Delta^0,i\Delta^5,i\Delta^6,i\Delta^7)$, which is why the relative phase $\pm i$ appears in
Eq.~(\ref{eq:groundstates}). Within the zero-mode sector the singlet $\Delta^0$ of the rotations
$SO(3)_T$ of the additional times and their triplet $\Delta^t$ are degenerate \emph{exactly}, to all
orders, because the hidden $SO(4)$ rotates one into the other. The symmetry breaking is
$SO(4)\to SO(3)_T$ together with $U(1)_F\to Z_2$, and there are four gapless modes. But the hidden
symmetry is a property of the zero-mode sector only. Every effect that involves a dependence of the
fields on $x0$, in the second stage of the program the excitations that travel around the hidden
ring, with energies of at least the inverse radius, breaks it down to $SO(3)_T$, under which the
hidden axis \emph{is} distinguished. Such effects lift the degeneracy between singlet and triplet and
give the three additional modes a mass. Which of the two then lies lower has not been determined;
that it is the singlet (\ref{eq:singlet}) is the content of hypothesis H1.
%%</1644>

\subsection{The space of the order parameter}
\label{sec:landscape}

The statements of Sec.~\ref{sec:ground} rest on the expansion (\ref{eq:GL}) near $T_c$. The CJT
potential allows them to be tested at zero temperature and away from small amplitudes. The program
builds the gap matrix (\ref{eq:gapmatrix}) from the Dirac matrices of each spacetime and evaluates the
potential for a constant attraction in each of the four channels, in units in which the gap of the
unitary state is $\Delta_0=1$ and $N(0)=1$. Table~\ref{tab:landscape} lists six configurations for
each spacetime, and Fig.~\ref{fig:landscape} shows a section through the space of the order
parameter.

\begin{table*}
\caption{\label{tab:landscape}The CJT potential at zero temperature on the eight-dimensional space of
the order parameter, from the Dirac matrices of both spacetimes; units $\Delta_0=N(0)=1$. For each
configuration: the two fourfold eigenvalues of $\MM\MM^\dagger$, the potential, the numbers of
vanishing and of positive eigenvalues of its matrix of second derivatives (with the largest
eigenvalue), and the real coordinates that span the flat directions.}
\begin{ruledtabular}
\footnotesize
\begin{tabular}{llcccl}
 & configuration & gaps$^2$ & $V$ & second derivatives & flat directions \\
\hline
{{ROWS:cjt_landscape_table}}
\end{tabular}
\end{ruledtabular}
\end{table*}

\begin{figure*}
\includegraphics[width=\textwidth]{fig_cjt_landscape}
\caption{\label{fig:landscape}A section through the space of the order parameter: directions
$(\cos\theta,\,e^{i\varphi}\sin\theta,\,0,\,0)$ on the extra axes $(x0,x5)$. (a), (b) The potential at
the amplitude that minimizes it in each direction, for the two spacetimes. The unitary states form
the flat valley at $-2$, which is the line of relative phase $\varphi=0$ on the $7+1$ spacetime and
$\varphi=90^\circ$ on the $4+4$ spacetime; the half-paired state is the saddle at $-1$. The two panels
differ only by the exchange $\varphi\leftrightarrow90^\circ-\varphi$. (c) The two fourfold gaps along
the section that contains the half-paired state.}
\end{figure*}

\emph{The unitary states are the minima}, with $V=-2$, that is, $-\nu N(0)\Delta_0^2/4$, whatever
their direction: the hidden-axis singlet, a state along $x5$ and rotated states have the same energy
to the accuracy of the calculation, on both spacetimes.

\emph{Four flat directions and four massive ones.} Around every unitary state four of the eight real
directions cost no energy: the phase, which belongs to the fermion number, and three rotations of the
direction. They are the directions of the four gapless collective modes. At the hidden-axis singlet
they are, on the $7+1$ spacetime, the imaginary part of $\Delta^0$ and the \emph{real} parts of
$\Delta^5,\Delta^6,\Delta^7$; on the $4+4$ spacetime, the imaginary part of $\Delta^0$ and the
\emph{imaginary} parts of the other three, because the hidden symmetry rotates $\Delta^0$ into
$i\Delta^t$. The other four directions, the amplitude and three that make the state non-unitary, have
equal curvature, {{GEN:cjt_massive_curvature}} $N(0)$.

\emph{A half-paired stationary state.} If $\eta_{ab}\Delta^a\Delta^b=0$, four of the eight internal
states are paired with the full gap and four are not paired at all, by Eq.~(\ref{eq:twogaps}). This
state is stationary and has half the condensation energy. It is unstable against opening the second
gap: the potential rises by {{GEN:cjt_half_up}} when the order parameter is multiplied by 1.02 or by
0.98, and falls by {{GEN:cjt_half_open}} when 0.05 is added to $\Delta^0$. It is the analog of the
non-unitary phases of superfluid $^3$He, which a magnetic field stabilizes \cite{Leggett1975}; here
the corresponding field would be a chemical potential for one of the generators of $SO(4)$.

\emph{The two spacetimes are mirror images.} In the direction
$(\cos\theta,e^{i\varphi}\sin\theta,0,0)$ the two gaps are, by Eq.~(\ref{eq:twogaps}), proportional to
$\sqrt{1\mp\sin2\theta\,|\sin\varphi|}$ on the $7+1$ spacetime and to
$\sqrt{1\mp\sin2\theta\,|\cos\varphi|}$ on the $4+4$ spacetime; the program compares these forms with
the eigenvalues computed from the matrices at every point of Fig.~\ref{fig:landscape}. A rotation of
the order parameter from the hidden axis towards $x5$ costs nothing if the two components have the
relative phase 0 on the $7+1$ spacetime, or $90^\circ$ on the $4+4$ spacetime. With any other
relative phase it costs energy, and with the opposite phase the path leads over the half-paired
state. The configuration that is unitary on one spacetime is far from unitary on the other (last row
of each block of Table~\ref{tab:landscape}).

The curvatures of the potential in the massive directions depend on how propagators away from the
stationary point are parametrized; they are not frequencies of collective modes. The number of flat
directions and the energies of stationary states do not depend on that choice.

\subsection{Collective modes and the sound of the superfluid}
\label{sec:sound}

The phase $\varphi$ of the order parameter is the Goldstone field of the broken fermion number. Its
excitation is a density wave, the first sound of the superfluid
\cite{Bogoliubov1947,Anderson1958}, with the velocity given by the compressibility,
$c_s^2=dP/d\rho$. In the weak-coupling regime the pressure is the degeneracy pressure of the fermions,
and in a dilute gas of bound pairs it is that of a weakly interacting Bose gas:
\begin{equation}
c_s^2=\frac{\kF^2}{3\,(m^2+\kF^2)}=\frac{v_F^2}{3},\qquad\text{or}\qquad c_s^2=\frac{4\pi a_{BB}\,n_B}{m_B^2},
\label{eq:sound}
\end{equation}
where $m_B$, $n_B$ and $a_{BB}>0$ are the mass, the density and the mutual scattering length of the
bound pairs. Equation (\ref{eq:sound}) is the input of the cosmological simulations of
Sec.~\ref{sec:dm}, in which the superfluid appears as a fluid with this sound speed. Besides the
sound mode both models have three gapless modes in which the direction of the order parameter in the
space of the extra axes oscillates. \SIG{On the $4+4$ spacetime they are protected by the hidden
symmetry only, and any dependence of the fields on $x0$ gives them a mass.}{On the $7+1$ spacetime
they are exact Goldstone modes of the flat model.}

\subsection{The field equations of the superfluid}
\label{sec:sfeqs}

For the reader who wants the complete set of equations in one place: in the unpaired phase the fields
obey Eqs.~(\ref{eq:fieldeqs}). In the paired phase, in the mean-field approximation of
Sec.~\ref{sec:gapeq}, the fermion field obeys the Bogoliubov-de Gennes equations
\begin{align}
i\partial_4\psi&=(\hh-\mu)\,\psi+\hat\Delta\,\psi^{\dagger T},\nonumber\\
\hat\Delta(\vec p)&=U(\vec p)^{*}\,\MM(\vec p)\,U(-\vec p)^{\dagger},
\label{eq:BdG}
\end{align}
in which $\hat\Delta$ is the $16\times16$ pair potential built from the gap matrix
(\ref{eq:gapmatrix}), and the gap obeys the self-consistency equation (\ref{eq:gap}). Near $T_c$ the
order parameter varies slowly and obeys the Ginzburg-Landau equation that follows from
Eqs.~(\ref{eq:GL}) and (\ref{eq:quartic}) by variation with respect to $\Delta^{a*}$,
\begin{align}
-K_{\rm GL}\,\nabla^2\Delta^a+a(T)\,\Delta^a+b\Big[&2\Big(\sum_c|\Delta^c|^2\Big)\Delta^a\nonumber\\
&-\big(\eta_{cd}\Delta^c\Delta^d\big)\,\eta_{ab}\,\Delta^{b*}\Big]=0 .
\label{eq:GLeq}
\end{align}
For a uniform unitary state it gives $a+b|\Delta|^2=0$, that is,
$|\Delta|^2=\frac{8\pi^2}{7\zeta(3)}T_c^2(1-T/T_c)$ near $T_c$.

\subsection{What has been established, and what has not}

Established, from the matrices of both spacetimes: the quartic invariant (\ref{eq:quartic}) with the
metric of the spacetime on the extra axes; the spectrum of gaps (\ref{eq:twogaps}) of a general
state; that the unitary states are the ground states and that the hidden-axis singlet is one of them;
four flat directions around every ground state; the half-paired stationary state.

%%<1671>
Not established: that the hidden axis is \emph{preferred}. On the flat $7+1$ spacetime the direction
of the order parameter is spontaneous. Also not determined are the frequencies and dampings of the
collective modes; only their number and the static energies were computed.
%%</1671>
%%<1644>
Not established: that the hidden axis is \emph{preferred}. In the zero-mode sector singlet and
triplet are exactly degenerate, and which is lower once the fields depend on $x0$ has not been
computed. Also not determined are the frequencies and dampings of the collective modes, and the
masses that the three additional gapless modes acquire when the fields depend on $x0$.
%%</1644>
